\documentclass[11pt,oneside]{article}

\usepackage[a4paper,margin=27mm,headheight=15pt]{geometry}
\usepackage[T1]{fontenc}
\usepackage[utf8]{inputenc}
\IfFileExists{libertinus.sty}{\usepackage{libertinus}}{\usepackage{lmodern}}
\usepackage{microtype}
\usepackage{amsmath,amssymb,amsthm,mathtools,mathrsfs}
\usepackage{bm}
\usepackage{booktabs,longtable,array,tabularx}
\usepackage{graphicx}
\usepackage{pgfplots}
\usepackage{xcolor}
\usepackage{enumitem}
\usepackage{fancyhdr}
\usepackage{titlesec}
\usepackage{listings}
\usepackage{xurl}
\usepackage{hyperref}
\usepackage[nameinlink,noabbrev]{cleveref}

\definecolor{linkblue}{RGB}{25,75,135}
\definecolor{shadegray}{RGB}{246,247,249}
\definecolor{rulegray}{RGB}{195,201,208}
\definecolor{deepgray}{RGB}{75,80,86}
\pgfplotsset{compat=1.18}

\hypersetup{
  colorlinks=true,
  linkcolor=linkblue,
  citecolor=linkblue,
  urlcolor=linkblue,
  pdfauthor={Vladimir Reshetnikov, ProveIt project},
  pdftitle={Dyadic Calculus for the Fabius Function: Connections and Alternative Representations}
}

\pagestyle{fancy}
\fancyhf{}
\fancyhead[L]{\small Fabius--Rvachev synthesis}
\fancyhead[R]{\small\nouppercase{\leftmark}}
\fancyfoot[C]{\thepage}
\renewcommand{\headrulewidth}{0.4pt}

\titleformat{\section}
  {\Large\bfseries}
  {\thesection}{0.7em}{}
  [\vspace{0.2ex}\color{rulegray}\titlerule]
\titleformat{\subsection}
  {\large\bfseries}
  {\thesubsection}{0.7em}{}
\titleformat{\subsubsection}
  {\normalsize\bfseries}
  {\thesubsubsection}{0.7em}{}
\titlespacing*{\section}{0pt}{2.8ex plus 1ex minus .2ex}{1.7ex plus .2ex}
\titlespacing*{\subsection}{0pt}{2.2ex plus .8ex minus .2ex}{1.0ex plus .2ex}

\setlist[itemize]{leftmargin=2em,itemsep=.25ex,topsep=.5ex}
\setlist[enumerate]{leftmargin=2.2em,itemsep=.25ex,topsep=.5ex}
\setlength{\parindent}{1.25em}
\setlength{\parskip}{0.15em}
\allowdisplaybreaks[3]

\newtheorem{theorem}{Theorem}[section]
\newtheorem{proposition}[theorem]{Proposition}
\newtheorem{lemma}[theorem]{Lemma}
\newtheorem{corollary}[theorem]{Corollary}
\theoremstyle{definition}
\newtheorem{definition}[theorem]{Definition}
\newtheorem{example}[theorem]{Example}
\theoremstyle{remark}
\newtheorem{remark}[theorem]{Remark}

\newenvironment{proofidea}
  {\par\smallskip\noindent\textit{Proof idea.}\ }
  {\hfill$\square$\par\smallskip}
\newenvironment{boxedremark}
  {\par\medskip\noindent\begin{center}\begin{minipage}{0.94\textwidth}\color{deepgray}\hrule\vspace{0.7ex}}
  {\vspace{0.7ex}\hrule\end{minipage}\end{center}\medskip}

\newcommand{\R}{\mathbb R}
\newcommand{\C}{\mathbb C}
\newcommand{\N}{\mathbb N}
\newcommand{\E}{\mathbb E}
\newcommand{\Prob}{\mathbb P}
\newcommand{\eps}{\varepsilon}
\newcommand{\dd}{\,\mathrm d}
\newcommand{\Up}{\operatorname{up}}
\newcommand{\sinc}{\operatorname{sinc}}
\newcommand{\wt}{\operatorname{wt}_2}
\newcommand{\fracpart}[1]{\left\{#1\right\}}
\newcommand{\poch}[3]{\left(#1;#2\right)_{#3}}
\newcommand{\qbinom}[3]{\genfrac{[}{]}{0pt}{}{#1}{#2}_{#3}}
\newcommand{\half}{\tfrac12}
\newcommand{\Bcal}{\mathcal B}
\newcommand{\Fglobal}{\mathcal F}
\newcommand{\Toperator}{\mathfrak T}
\newcommand{\Qoperator}{\mathfrak Q}
\newcommand{\coeff}[2]{\left[#1^{#2}\right]}

\title{\bfseries Dyadic Calculus for the Fabius Function\\[0.35em]
\Large Connections among exact values, Appell blocks, Thue--Morse differences, Gaussian interpolation, and discrete limits}
\author{A complementary article to \emph{The Fabius Function and Rvachev's up-Function}}
\date{25 August 2026}

\begin{document}
\maketitle

\begin{abstract}
The exact formulae for the values of the Fabius function at dyadic rationals,
the binary-reduction series at an arbitrary argument, and the discrete-limit
formula involving Gaussian binomial coefficients look at first like three
rather different constructions.  They are, in fact, organized by one finite
polynomial identity and one regular summability argument.

The finite identity is most transparent after introducing the Appell moment
polynomials
\[
  A_n(\theta)=\E (X+\theta)^n,
\]
where $X$ has the Fabius distribution.  The Taylor-reduction polynomial at
scale $n$ is exactly $2^{-\binom n2}A_n(2^n y)/n!$.  Consequently, the global
binary expansion becomes a rapidly convergent, digit-sparse Appell series;
for a dyadic rational it terminates, so the finite dyadic formula is literally
the terminating case of the arbitrary-argument expansion.  Replacing each
Appell polynomial by its finite Thue--Morse/Gaussian-binomial representation
gives the nested $q$-series term by term, not merely after taking a limit.

The discrete-limit formula uses the same finite Gaussian stencil in a different
way.  After an exact reindexing it is a signed Toeplitz average of shifted
uniform-spline approximants.  Its row-generating polynomial is
\[
  \frac{\poch{z/2}{1/2}{n}}{\poch{1/2}{1/2}{n}},
\]
which simultaneously explains unit row mass, bounded total variation, and the
appearance of $q$-Pochhammer symbols.  This article develops these connections,
derives quantitative error bounds, and adds several alternative exact
representations: a finite composition sum obtained directly from the moment
product, a lower-Hessenberg determinant, a Bernoulli--Bell-polynomial formula,
and Fourier and Laplace inversion formulae.
\end{abstract}

\tableofcontents

\section{Orientation and notation}

\subsection{What is being connected}

The companion article and the Lean development establish, among many other
facts, the following three families of formulae.

\begin{enumerate}
\item Exact rational expressions for $F(m/2^n)$, including formulae in which a
  Gaussian binomial coefficient at base $1/2$ multiplies a finite signed
  Thue--Morse power sum.
\item An exact infinite expansion for $F(x)$, or for the signed
  global extension $\Fglobal(x)$, obtained by repeatedly reducing along the
  binary expansion of $x$.  Substituting the finite $q$-binomial expression
  into every reduction block produces a nested infinite $q$-series.
\item A discrete limit in which a finite $q$-binomial sum of increasingly fine
  Thue--Morse splines converges to $\Fglobal(x)$ at an arbitrary real argument.
\end{enumerate}

The purpose here is not to repeat the proofs already presented there.  We
instead isolate the common algebraic mechanism, compress the formulae into a
small collection of reusable objects, and then unfold those objects in several
new directions.

Throughout,
\[
  \rho=\frac12,
  \qquad
  \poch{a}{\rho}{n}=\prod_{j=0}^{n-1}(1-a\rho^j),
  \qquad
  \qbinom{n}{k}{\rho}
  =\frac{\poch{\rho}{\rho}{n}}
  {\poch{\rho}{\rho}{k}\poch{\rho}{\rho}{n-k}}.
\]
The empty product is $1$.  We write
\[
  \eps_r=(-1)^{\wt(r)},
\]
where $\wt(r)$ is the number of $1$'s in the binary expansion of $r$.
Thus $(\eps_r)_{r\ge0}$ is the $\{\pm1\}$ Thue--Morse sequence.


\subsection{The half-base notation is Mersenne arithmetic}

At the special base $\rho=1/2$, the $q$-Pochhammer symbol is simply a compact
normalization of a product of Mersenne numbers:
\begin{equation}\label{eq:Mersenne-Pochhammer}
 \poch{1/2}{1/2}{n}
 =\prod_{j=1}^{n}(1-2^{-j})
 =2^{-\binom{n+1}{2}}\prod_{j=1}^{n}(2^j-1).
\end{equation}
In particular,
\begin{equation}\label{eq:main-denom-Mersenne}
 2^{n^2}\poch{1/2}{1/2}{n}
 =2^{\binom n2}\prod_{j=1}^{n}(2^j-1).
\end{equation}
The factors $2^j-1$ appearing separately in moment recurrences and
composition formulae are therefore the same arithmetic content that the
$q$-Pochhammer symbol packages in the closed Gaussian-binomial formula.
Likewise,
\begin{equation}\label{eq:qbinomial-inversion}
 \qbinom{n}{k}{1/2}
 =2^{-k(n-k)}\qbinom{n}{k}{2}.
\end{equation}
Thus the half-base Gaussian coefficients are ordinary integral Gaussian
coefficients at base $2$, accompanied by an explicit power of $2$.  The
$q$-notation is not introducing mysterious irrational constants; it records
the geometry of the dyadic interpolation grid and keeps its Mersenne
normalizations factored.

\subsection{The probabilistic normalization}

Let $U_1,U_2,\ldots$ be independent random variables, each uniformly
distributed on $[0,1]$, and put
\begin{equation}\label{eq:X-definition}
  X=\sum_{j=1}^{\infty}\frac{U_j}{2^j}.
\end{equation}
Then $0\le X\le1$, and the distribution function of $X$ is the Fabius
function:
\[
  F(x)=\Prob(X\le x).
\]
This representation fixes all normalizations used below.  Define the moments
\begin{equation}\label{eq:dn-definition}
  d_n=\E X^n,
  \qquad d_0=1,
\end{equation}
and their exponential generating function
\begin{equation}\label{eq:G-definition}
  G(z)=\E e^{zX}=\sum_{n=0}^{\infty}d_n\frac{z^n}{n!}.
\end{equation}
Independence in \eqref{eq:X-definition} gives the entire product
\begin{equation}\label{eq:G-product}
  G(z)=\prod_{j=1}^{\infty}E\!\left(\frac{z}{2^j}\right),
  \qquad
  E(w)=\frac{e^w-1}{w}.
\end{equation}
Shifting the product by one factor yields
\begin{equation}\label{eq:G-dilation}
  G(2z)=E(z)G(z).
\end{equation}
Comparing coefficients gives the basic moment recurrence
\begin{equation}\label{eq:moment-recurrence}
  (2^n-1)\frac{d_n}{n!}
  =\sum_{k=0}^{n-1}\frac{d_k}{k!(n-k+1)!}
  \qquad(n\ge1).
\end{equation}

One of the fundamental exact dyadic identities is
\begin{equation}\label{eq:inverse-dyadic-moment}
  F(2^{-n})=2^{-\binom n2}\frac{d_n}{n!}.
\end{equation}
It will reappear in nearly every section, but in different disguises.

\begin{boxedremark}
The product \eqref{eq:G-product}, the recurrence
\eqref{eq:moment-recurrence}, the dyadic identity
\eqref{eq:inverse-dyadic-moment}, the finite $q$-binomial formulae, and the
spline limits are all formally proved in the cited Lean development.  Some of
the regroupings below---notably the Appell compression, determinant form, and
explicit error estimate for the Toeplitz limit---are elementary deductions
from those formalized identities and are presented here in human-readable
form.
\end{boxedremark}

\section{The master Taylor block}

\subsection{The reduction polynomial}

For $n\ge0$, define
\begin{equation}\label{eq:Rn-definition}
 R_n(y)=
 \sum_{k=0}^{n}
 2^{\binom{k+1}{2}-\binom{n-k}{2}}
 \frac{d_{n-k}}{(n-k)!}\frac{y^k}{k!}.
\end{equation}
This is the polynomial that occurs in the exact Taylor reduction.  On the
appropriate dyadic shell, if $y=x-2^{-n}$, the global Fabius extension satisfies
an identity of the form
\[
  \Fglobal(x)=-\Fglobal(y)+R_n(y).
\]
Only the polynomial itself will be needed here.

Putting $y=0$ in \eqref{eq:Rn-definition} gives
\[
  R_n(0)=2^{-\binom n2}\frac{d_n}{n!}=F(2^{-n}),
\]
so every inverse-dyadic value is the constant term of one reduction block.
The less obvious fact is that the whole block has an equally simple meaning.

\subsection{Appell compression}

\begin{definition}[Moment Appell polynomials]\label{def:Appell}
For $n\ge0$ and $\theta\in\C$, let
\begin{equation}\label{eq:An-definition}
  A_n(\theta)=\E(X+\theta)^n
  =\sum_{k=0}^{n}\binom nk d_{n-k}\theta^k.
\end{equation}
Equivalently,
\begin{equation}\label{eq:Appell-egf}
  \sum_{n=0}^{\infty}A_n(\theta)\frac{z^n}{n!}
  =e^{\theta z}G(z).
\end{equation}
\end{definition}

The family is Appell in the classical sense:
\begin{equation}\label{eq:Appell-derivative}
  A_n'(\theta)=nA_{n-1}(\theta),
  \qquad
  A_n(\theta+c)=\sum_{j=0}^{n}\binom nj A_{n-j}(\theta)c^j.
\end{equation}
Since $X$ and $1-X$ have the same distribution,
\begin{equation}\label{eq:Appell-reflection}
  A_n(\theta)=(-1)^nA_n(-1-\theta).
\end{equation}

\begin{theorem}[Appell compression of the reduction block]\label{thm:Appell-compression}
For every $n\ge0$ and $\theta\in\C$,
\begin{equation}\label{eq:Appell-compression}
  R_n(2^{-n}\theta)
  =2^{-\binom n2}\frac{A_n(\theta)}{n!}.
\end{equation}
Equivalently, for arbitrary $y$,
\begin{equation}\label{eq:R-Appell-general}
  R_n(y)=2^{-\binom n2}\frac{A_n(2^ny)}{n!}.
\end{equation}
\end{theorem}

\begin{proof}
Substitute $y=2^{-n}\theta$ into \eqref{eq:Rn-definition}.  The power of $2$
in the $k$th term becomes
\[
 \binom{k+1}{2}-\binom{n-k}{2}-nk=-\binom n2,
\]
independently of $k$.  Factoring out $2^{-\binom n2}$ leaves
\[
 \sum_{k=0}^{n}\frac{d_{n-k}\theta^k}{(n-k)!k!}
 =\frac1{n!}\sum_{k=0}^{n}\binom nkd_{n-k}\theta^k
 =\frac{A_n(\theta)}{n!}.
\]
\end{proof}

This elementary cancellation is the main compression used throughout the
article.  A triangular-looking polynomial with apparently irregular powers of
$2$ is simply a rescaled shifted moment.

\subsection{Immediate consequences}

For $0\le\theta<1$, positivity and bounded support give
\begin{equation}\label{eq:Appell-bound}
  0\le A_n(\theta)\le 2^n.
\end{equation}
Consequently,
\begin{equation}\label{eq:R-superbound}
  0\le R_n(2^{-n}\theta)
  \le \frac{2^{n-\binom n2}}{n!}.
\end{equation}
This estimate is much sharper than the geometric estimate normally used to
prove convergence of the binary-reduction series.

The first few Appell polynomials are
\begin{align*}
 A_0(\theta)&=1,\\
 A_1(\theta)&=\theta+\frac12,\\
 A_2(\theta)&=\theta^2+\theta+\frac5{18},\\
 A_3(\theta)&=\theta^3+\frac32\theta^2+\frac56\theta+\frac16,\\
 A_4(\theta)&=\theta^4+2\theta^3+\frac53\theta^2
              +\frac23\theta+\frac{143}{1350}.
\end{align*}
They encode all Taylor blocks of the corresponding orders by a single change
of scale.

\section{The binary expansion as a sparse Appell series}

\subsection{Prefixes, bits, and tails}

Fix $x\in[0,1)$.  Use the canonical binary expansion that is not eventually
all $1$'s,
\begin{equation}\label{eq:binary-x}
  x=0.b_1b_2b_3\ldots{}_2,
  \qquad b_m\in\{0,1\}.
\end{equation}
Let
\begin{equation}\label{eq:prefix-tail}
  a_m=\lfloor2^m x\rfloor,
  \qquad
  \theta_m=\fracpart{2^m x},
  \qquad
  y_m=x-\frac{a_m}{2^m}=2^{-m}\theta_m.
\end{equation}
Then
\begin{equation}\label{eq:bit-prefix-relation}
  a_m=2a_{m-1}+b_m,
  \qquad
  \theta_m=0.b_{m+1}b_{m+2}\ldots{}_2.
\end{equation}

The binary-reduction series in the formal development has coefficient
\begin{equation}\label{eq:source-coefficient}
  c_m(x)=\eps_{a_m}\bigl(2a_{m-1}-a_m\bigr)
  \qquad(m\ge1).
\end{equation}
Using \eqref{eq:bit-prefix-relation}, this simplifies completely:
\begin{equation}\label{eq:coefficient-simplified}
  c_m(x)=
  \begin{cases}
   0,&b_m=0,\\
   \eps_{a_{m-1}},&b_m=1.
  \end{cases}
\end{equation}
Indeed, when $b_m=1$ one has
$\eps_{a_m}=-\eps_{a_{m-1}}$ and $2a_{m-1}-a_m=-1$.
Thus only the positions of the $1$-bits contribute.

\subsection{The digit--Appell formula}

\begin{theorem}[Digit-sparse Appell expansion]\label{thm:digit-Appell}
For every $x\in[0,1)$,
\begin{equation}\label{eq:digit-Appell}
 \boxed{
 F(x)=
 \sum_{m=1}^{\infty}
 b_m\,(-1)^{b_1+\cdots+b_{m-1}}
 \frac{2^{-\binom m2}}{m!}
 A_m(\theta_m).}
\end{equation}
The series converges absolutely, and in fact uniformly with respect to
$x\in[0,1)$.
\end{theorem}

\begin{proof}
The exact binary-reduction series is
\[
 F(x)=\sum_{m\ge1}c_m(x)R_m(y_m).
\]
Insert \eqref{eq:coefficient-simplified}, use
$\eps_{a_{m-1}}=(-1)^{b_1+\cdots+b_{m-1}}$, and then apply
\cref{thm:Appell-compression} with $y_m=2^{-m}\theta_m$.
Absolute and uniform convergence follow from
\eqref{eq:Appell-bound} and the Weierstrass test.
\end{proof}

Formula \eqref{eq:digit-Appell} is a compact arbitrary-argument representation
that contains no $q$-notation at all.  The $q$-series will emerge by replacing
each $A_m$ by an exactly equal finite Gaussian-binomial expression.

\subsection{A superexponential remainder bound}

Let $S_N(x)$ denote the truncation of \eqref{eq:digit-Appell} after $m=N$.
From \eqref{eq:Appell-bound},
\[
 |F(x)-S_N(x)|
 \le\sum_{m=N+1}^{\infty}\frac{2^{m-\binom m2}}{m!}.
\]
The ratio of consecutive majorants is
\[
 \frac{2^{m+1-\binom{m+1}{2}}/(m+1)!}
      {2^{m-\binom m2}/m!}
 =\frac{2^{1-m}}{m+1}\le\frac12
 \qquad(m\ge1).
\]
Therefore, for $N\ge0$,
\begin{equation}\label{eq:digit-tail-bound}
 \boxed{
  |F(x)-S_N(x)|
  \le 2\,\frac{2^{N+1-\binom{N+1}{2}}}{(N+1)!}.}
\end{equation}
The right-hand side decays faster than geometrically in $N$.  In practical
computation, the binary digit series is therefore far more rapidly convergent
than the elementary residual estimate $O(2^{-N})$ suggests.

\subsection{Dyadic rationals are the terminating case}

Suppose
\[
 x=\frac{M}{2^N}=0.b_1b_2\ldots b_N000\ldots{}_2,
 \qquad 0\le M<2^N.
\]
Then $b_m=0$ for $m>N$, and \eqref{eq:digit-Appell} terminates:
\begin{equation}\label{eq:dyadic-digit-Appell}
 \boxed{
 F\!\left(\frac{M}{2^N}\right)
 =\sum_{\substack{1\le m\le N\\b_m=1}}
 (-1)^{b_1+\cdots+b_{m-1}}
 \frac{2^{-\binom m2}}{m!}
 A_m\!\left(0.b_{m+1}\ldots b_N{}_2\right).}
\end{equation}
Thus the exact bit-recursive evaluator, the finite exact dyadic formula, and
the arbitrary-argument series are not separate phenomena: they are the finite,
recursive, and infinite readings of the same identity.

\begin{example}[The value at $5/16$]\label{ex:five-sixteen}
The terminating binary expansion is
\[
 \frac5{16}=0.0101_2.
\]
Only $m=2$ and $m=4$ contribute.  At $m=2$ the previous prefix has even
weight and $\theta_2=1/4$; at $m=4$ the previous prefix is $010_2$, so the
sign is negative and $\theta_4=0$.  Hence
\begin{align*}
 F\!\left(\frac5{16}\right)
 &=\frac{2^{-1}}{2!}A_2\!\left(\frac14\right)
   -\frac{2^{-6}}{4!}A_4(0)\\
 &=\frac14\left(\frac5{18}+\frac14+\frac1{16}\right)
   -\frac1{1536}\frac{143}{1350}\\
 &=\frac{305857}{2073600}.
\end{align*}
The first contribution accounts for the prefix $01$ together with the
remaining tail $01$; the second contribution corrects for the final $1$-bit.
\end{example}

\section{Why Thue--Morse and Gaussian binomial coefficients appear}

Two independent annihilation mechanisms are superposed in the nested exact
formula.  The inner Thue--Morse sum is an additive finite-difference operator;
the outer Gaussian-binomial sum is a divided-difference operator on a geometric
grid.  Each kills lower-degree terms.  Their composition isolates precisely
the coefficient needed by the Taylor block.

\subsection{The Thue--Morse difference operator}

For a polynomial or entire function $P$, define
\begin{equation}\label{eq:TM-operator}
  (\Toperator_kP)(c)
  =\sum_{r=0}^{2^k-1}\eps_rP(c+r).
\end{equation}
Writing $D=\frac{\dd}{\dd c}$ and using the binary digits of $r$ gives the
operator factorization
\begin{equation}\label{eq:TM-factorization}
  \Toperator_k
  =\prod_{j=0}^{k-1}\left(1-e^{2^jD}\right).
\end{equation}
Indeed,
\[
 \sum_{r<2^k}\eps_re^{rz}
 =\prod_{j=0}^{k-1}(1-e^{2^jz}).
\]
Since every factor in \eqref{eq:TM-factorization} begins with
$-2^jD$, the operator annihilates all polynomials of degree less than $k$ and
satisfies
\begin{equation}\label{eq:first-TM-moment}
  \sum_{r=0}^{2^k-1}\eps_rr^k
  =(-1)^k2^{\binom k2}k!.
\end{equation}
More generally, for fixed $n$,
\begin{equation}\label{eq:TM-degree-n}
  c\longmapsto
  \frac1{(n+k)!}
  \sum_{r<2^k}\eps_r(r+c)^{n+k}
\end{equation}
has degree at most $n$, and its leading coefficient is
\begin{equation}\label{eq:TM-leading-coeff}
  \coeff{c}{n}\left(
  \frac1{(n+k)!}\sum_{r<2^k}\eps_r(r+c)^{n+k}
  \right)
  =\frac{(-1)^k2^{\binom k2}}{n!}.
\end{equation}

A symmetric form of the operator is also useful.  With
$\operatorname{sinhc}(z)=\sinh(z)/z$,
\begin{equation}\label{eq:TM-sinhc}
 \Toperator_k
 =(-1)^k2^{\binom k2}D^k
 e^{(2^k-1)D/2}
 \prod_{j=0}^{k-1}
 \operatorname{sinhc}(2^{j-1}D).
\end{equation}
This shows that the Thue--Morse block is a $k$th derivative, followed by an
even differential correction and evaluation about the midpoint of the block.
It also explains why centered and translated versions of the formula are so
natural.

\subsection{Geometric divided differences}

Let $0<q<1$ for a moment, and consider the geometric nodes
\[
  1,q^{-1},q^{-2},\ldots,q^{-n}.
\]
The leading coefficient of any polynomial $P$ of degree at most $n$ is its
$n$th divided difference on these nodes.  Lagrange interpolation gives
\begin{equation}\label{eq:q-leading-extractor-general}
 \coeff{t}{n}P(t)
 =\frac{(-1)^nq^{\binom{n+1}{2}}}{\poch{q}{q}{n}}
  \sum_{k=0}^{n}
  (-1)^kq^{\binom k2}\qbinom{n}{k}{q}P(q^{-k}).
\end{equation}
For completeness, the barycentric denominator is
\begin{equation}\label{eq:barycentric-denominator}
 \frac1{\displaystyle\prod_{j\ne k}(q^{-k}-q^{-j})}
 =\frac{(-1)^{n-k}q^{\binom{n+1}{2}+\binom k2}}
 {\poch{q}{q}{n}}\qbinom{n}{k}{q}.
\end{equation}
At $q=1/2$, the nodes are $1,2,4,\ldots,2^n$ and the interpolation weights are
exactly Gaussian binomial coefficients at base $1/2$, decorated by triangular
powers of $2$.

The same identity follows immediately from the finite $q$-binomial theorem
\begin{equation}\label{eq:finite-q-binomial}
 \sum_{k=0}^{n}(-1)^kq^{\binom k2}\qbinom{n}{k}{q}z^k
 =\poch{z}{q}{n}.
\end{equation}
If $z=q^{-m}$, the right-hand side is zero for $m<n$, while
\begin{equation}\label{eq:q-survivor}
  \poch{q^{-n}}{q}{n}
  =(-1)^nq^{-\binom{n+1}{2}}\poch{q}{q}{n}.
\end{equation}
Thus the stencil kills every lower monomial and normalizes the top one.

\subsection{The two annihilators together}

The inner sum \eqref{eq:TM-degree-n} first lowers degree by $k$.  Its remaining
dependence is a polynomial of degree at most $n$.  The outer geometric stencil
\eqref{eq:q-leading-extractor-general} then kills all but one prescribed
coefficient of that degree-$n$ polynomial.  Schematically,
\[
 \boxed{
 \text{power of degree }n+k
 \xrightarrow{\ \Toperator_k\ }
 \text{polynomial of degree }n
 \xrightarrow{\ \Qoperator_n\ }
 \text{one coefficient}.}
\]
This is the structural reason that a formula about a smooth distribution
function involves both Thue--Morse signs and Gaussian binomial coefficients.
Neither ingredient is ornamental: each is the coefficient vector of a finite
annihilator.

Translation invariance is now transparent.  Translating the argument of the
intermediate degree-$n$ polynomial changes only lower coefficients after the
target coefficient is fixed.  The geometric extractor annihilates precisely
those changes.  The Lean proof strengthens this observation by treating the
translation as an indeterminate and proving that the complete numerator
polynomial is constant.

\section{The finite $q$-formula as an Appell identity}

\subsection{A centered master formula}

The exact inverse-dyadic formula is most revealing when one retains the Taylor
variable instead of setting it to zero.

\begin{theorem}[Centered Gaussian-binomial representation of $A_n$]
\label{thm:centered-q-Appell}
For every $n\ge0$ and $\theta\in\C$,
\begin{equation}\label{eq:centered-q-Appell}
\boxed{
 \frac{A_n(\theta)}{n!}
 =\frac{2^{-\binom{n+1}{2}}}{\poch{1/2}{1/2}{n}}
 \sum_{k=0}^{n}
 \frac{\qbinom{n}{k}{1/2}}
 {4^{\binom k2}(n+k)!}
 \sum_{r=0}^{2^k-1}
 \eps_r\bigl(r-2^k(1+\theta)\bigr)^{n+k}.}
\end{equation}
\end{theorem}

\begin{proof}
The polynomial version of the formalized $q$-binomial Taylor identity is
\[
 R_n(y)=
 \frac1{2^{n^2}\poch{1/2}{1/2}{n}}
 \sum_{k=0}^{n}
 \frac{\qbinom{n}{k}{1/2}}
 {4^{\binom k2}(n+k)!}
 \sum_{r<2^k}\eps_r
 \bigl(r-2^k-2^{n+k}y\bigr)^{n+k}.
\]
Set $y=2^{-n}\theta$ and use
\cref{thm:Appell-compression}.  Multiplying by $2^{\binom n2}$ changes the
outer power $2^{-n^2}$ into $2^{-\binom{n+1}{2}}$, while the inner shift becomes
$-2^k(1+\theta)$.
\end{proof}

At $\theta=0$, multiplying \eqref{eq:centered-q-Appell} by
$2^{-\binom n2}$ gives the exact centered inverse-dyadic formula
\begin{equation}\label{eq:centered-inverse-dyadic}
 \boxed{
 F(2^{-n})
 =\frac1{2^{n^2}\poch{1/2}{1/2}{n}}
 \sum_{k=0}^{n}
 \frac{\qbinom{n}{k}{1/2}}
 {4^{\binom k2}(n+k)!}
 \sum_{r=0}^{2^k-1}
 \eps_r(r-2^k)^{n+k}.}
\end{equation}
Thus the formula for one special value is merely the constant-term
specialization of a formula for the entire Appell polynomial.

\subsection{The arbitrary-translation form}

The companion development proves an even stronger polynomial identity.  A
common translation parameter may be inserted into all the raw Thue--Morse
blocks without changing the result.

\begin{theorem}[Raw translated representation]\label{thm:raw-q-Appell}
For every $n\ge0$ and every $\theta,\eta\in\C$,
\begin{equation}\label{eq:raw-q-Appell}
\boxed{
 \frac{A_n(\theta)}{n!}
 =\frac{(-1)^n2^{-\binom{n+1}{2}}}{\poch{1/2}{1/2}{n}}
 \sum_{k=0}^{n}
 \frac{\qbinom{n}{k}{1/2}}
 {4^{\binom k2}(n+k)!}
 \sum_{r=0}^{2^k-1}
 \eps_r\bigl(r+\eta+2^k\theta\bigr)^{n+k}.}
\end{equation}
The right-hand side is independent of $\eta$.
\end{theorem}

Putting $\theta=0$ gives
\begin{equation}\label{eq:raw-inverse-dyadic}
 F(2^{-n})
 =\frac1{(-2)^{n^2}\poch{1/2}{1/2}{n}}
 \sum_{k=0}^{n}
 \frac{\qbinom{n}{k}{1/2}}
 {4^{\binom k2}(n+k)!}
 \sum_{r=0}^{2^k-1}\eps_r(r+\eta)^{n+k}.
\end{equation}
The identity $(-1)^n=(-1)^{n^2}$ accounts for the denominator
$(-2)^{n^2}$.

\begin{remark}[A scale-dependent translation]
Because \eqref{eq:raw-q-Appell} is an identity separately for each $n$, one
may choose a different common translation $\eta_n$ at every outer scale in an
infinite expansion.  This does not change a single term.  In exact arithmetic
this is merely a gauge freedom; in floating-point arithmetic it may be used to
reduce the size of intermediate powers.  The translation must remain common
to the inner $k$-sum for a fixed $n$.
\end{remark}


\section{One-shot and digitwise formulae for a general dyadic rational}

The formula \eqref{eq:centered-inverse-dyadic} treats the numerator $1$.
There is also a one-shot formula for a general numerator.  For
$0\le M<2^N$,
\begin{equation}\label{eq:one-shot-general-dyadic}
\boxed{
 F\!\left(\frac{M}{2^N}\right)
 =\frac1{2^{N^2}\poch{1/2}{1/2}{N}}
 \sum_{k=0}^{N}
 \frac{\qbinom{N}{k}{1/2}}
 {4^{\binom k2}(N+k)!}
 \sum_{r=0}^{M2^k-1}
 \eps_r(r-M2^k)^{N+k}.}
\end{equation}
A common translation may be inserted in all the inner powers; the complete
outer numerator is unchanged.  Formula \eqref{eq:one-shot-general-dyadic}
uses the denominator exponent $N$ once and for all.  In contrast,
\eqref{eq:dyadic-digit-Appell} uses one lower-order block for each $1$-bit of
$M$.  The two organizations are related by the canonical dyadic decomposition
of the prefix interval.

Write
\[
 \frac{M}{2^N}=0.b_1\ldots b_N{}_2,
 \qquad
 a_m=\lfloor2^mM/2^N\rfloor.
\]
For every function $P$ on the nonnegative integers and every $k\ge0$,
\begin{align}\label{eq:prefix-block-decomposition}
 \sum_{r=0}^{M2^k-1}\eps_rP(r)
 &=\sum_{\substack{1\le m\le N\\b_m=1}}
 \eps_{a_{m-1}}
 \sum_{s=0}^{2^{N-m+k}-1}\eps_s
 P\!\left(a_{m-1}2^{N-m+k+1}+s\right).
\end{align}
Indeed, the interval $[0,M)$ is the disjoint union, over its $1$-bits, of
\[
 \left[a_{m-1}2^{N-m+1},
       a_{m-1}2^{N-m+1}+2^{N-m}\right),
 \qquad b_m=1.
\]
After multiplication by $2^k$, every block starts at a multiple of twice its
length.  Hence adding the local coordinate $s$ creates no binary carry and
\[
 \eps_{a_{m-1}2^{N-m+k+1}+s}
 =\eps_{a_{m-1}}\eps_s.
\]
This proves \eqref{eq:prefix-block-decomposition}.

Substituting \eqref{eq:prefix-block-decomposition} into the one-shot formula
splits every long Thue--Morse prefix into complete blocks carrying exactly the
signs $\eps_{a_{m-1}}$ that occur in the digit formula.  The remaining outer
Gaussian stencil is still of order $N$; applying the polynomial Taylor
identity to each rescaled block lowers it to its natural bit scale $m$ and
produces \eqref{eq:dyadic-digit-Appell}.  Thus:
\begin{itemize}
\item \eqref{eq:one-shot-general-dyadic} groups the computation by polynomial
  degree $N+k$;
\item \eqref{eq:dyadic-digit-Appell} groups it by the occupied binary scales of
  the numerator;
\item \eqref{eq:prefix-block-decomposition} is the combinatorial bridge between
  the two groupings.
\end{itemize}

For $M=5$ and $N=4$, the one-shot expression has one outer sum of order $4$
and inner prefixes of lengths $5\cdot2^k$.  The digitwise expression has only
two active blocks, at scales $m=2$ and $m=4$, as shown in
\cref{ex:five-sixteen}.  Their equality is a nontrivial finite rearrangement,
not a limiting assertion.

\section{From exact dyadic values to the arbitrary-argument $q$-series}

The connection can now be made without any further analysis: substitute
\cref{thm:centered-q-Appell} or \cref{thm:raw-q-Appell} into the digit--Appell
series term by term.

\subsection{The centered nested series}

\begin{theorem}[Centered digit--$q$ expansion]\label{thm:centered-digit-q}
For $x\in[0,1)$, with notation \eqref{eq:binary-x}--\eqref{eq:prefix-tail},
\begin{align}\label{eq:centered-digit-q}
 F(x)
 &=\sum_{m=1}^{\infty}
 b_m\eps_{a_{m-1}}
 \frac1{2^{m^2}\poch{1/2}{1/2}{m}}
 \sum_{k=0}^{m}
 \frac{\qbinom{m}{k}{1/2}}
 {4^{\binom k2}(m+k)!}
 \notag\\[-0.25em]
 &\hspace{5.6em}\times
 \sum_{r=0}^{2^k-1}
 \eps_r\bigl(r-2^k(1+\theta_m)\bigr)^{m+k}.
\end{align}
The equality holds term by term with \eqref{eq:digit-Appell}.
\end{theorem}

\begin{proof}
In the $m$th term of \eqref{eq:digit-Appell}, multiply the prefactor
$2^{-\binom m2}$ by the prefactor
$2^{-\binom{m+1}{2}}$ in \eqref{eq:centered-q-Appell}.  Their product is
$2^{-m^2}$.  Everything else is a direct substitution.
\end{proof}

If $x=M/2^N$ has a terminating binary expansion, then
\eqref{eq:centered-digit-q} terminates after $m=N$.  The ``infinite'' global
$q$-series is therefore exactly a finite dyadic formula whenever its argument
is dyadic.  Conversely, allowing infinitely many nonzero binary digits turns
the same finite block identity into the arbitrary-argument series.

\subsection{The raw translated nested series}

For any fixed $\eta\in\C$, the raw form is
\begin{align}\label{eq:raw-digit-q}
 F(x)
 &=\sum_{m=1}^{\infty}
 b_m\eps_{a_{m-1}}
 \frac1{(-2)^{m^2}\poch{1/2}{1/2}{m}}
 \sum_{k=0}^{m}
 \frac{\qbinom{m}{k}{1/2}}
 {4^{\binom k2}(m+k)!}
 \notag\\[-0.25em]
 &\hspace{5.6em}\times
 \sum_{r=0}^{2^k-1}
 \eps_r\bigl(r+\eta+2^k\theta_m\bigr)^{m+k}.
\end{align}
Again, every individual $m$th summand equals the corresponding Appell summand
in \eqref{eq:digit-Appell}; no interchange of infinite sums is needed to prove
the identity.

\subsection{A useful commutative diagram}

The relations may be summarized as follows:
\[
\begin{array}{ccc}
 R_m(2^{-m}\theta)
 &\xleftrightarrow{\quad\text{Appell compression}\quad}&
 2^{-\binom m2}A_m(\theta)/m!
 \\[0.8em]
 \Big\updownarrow\scriptstyle\text{finite Gaussian/Thue--Morse identity}
 &&
 \Big\downarrow\scriptstyle\text{insert binary tail }\theta=\theta_m
 \\[0.8em]
 \text{finite nested }q\text{-block}
 &\xrightarrow{\quad\text{sum over the 1-bits}\quad}&
 F(x).
\end{array}
\]
Setting $\theta=0$ in the top row gives $F(2^{-m})$.  Taking a terminating
binary expansion makes the right-hand summation finite.  Taking an arbitrary
binary expansion leaves an absolutely convergent infinite series.  These are
specializations, not analogies.

\section{Alternative exact representations of inverse-dyadic values}

The $q$-binomial formula is only one exact representation of
$F(2^{-n})$.  The moment product \eqref{eq:G-product} leads to several others,
each exposing a different combinatorial structure.

\subsection{Coefficient extraction from the moment product}

Let
\begin{equation}\label{eq:an-definition}
  a_n=\frac{d_n}{n!}=2^{\binom n2}F(2^{-n}).
\end{equation}
Then
\begin{equation}\label{eq:an-coeff}
  a_n=\coeff{z}{n}\prod_{j=1}^{\infty}
  \left(\sum_{r=0}^{\infty}
  \frac{z^r}{2^{jr}(r+1)!}\right).
\end{equation}
Expanding by the set of scales at which a nonconstant term is chosen gives
\begin{align}\label{eq:scale-selection}
 a_n
 &=\sum_{\ell=1}^{n}
   \sum_{\substack{r_1+\cdots+r_\ell=n\\r_i\ge1}}
   \frac1{\prod_{i=1}^{\ell}(r_i+1)!}
   \sum_{1\le j_1<\cdots<j_\ell}
   2^{-\sum_{i=1}^{\ell}j_ir_i},
 \qquad(n\ge1).
\end{align}
The infinite inner sum is an elementary geometric simplex sum.  If
$t_h=r_h+\cdots+r_\ell$, then
\begin{equation}\label{eq:geometric-simplex}
 \sum_{1\le j_1<\cdots<j_\ell}
 2^{-\sum_i j_ir_i}
 =\prod_{h=1}^{\ell}\frac1{2^{t_h}-1}.
\end{equation}
To see this, put
$j_i=i+g_1+\cdots+g_i$ with $g_h\ge0$; the resulting independent geometric
series contribute $(1-2^{-t_h})^{-1}$, and their powers of $2$ cancel the
initial offset.

Reversing the composition $(r_1,\ldots,r_\ell)$ changes tail sums into prefix
sums.  We obtain the finite composition formula
\begin{theorem}[Composition representation]\label{thm:composition}
Let $\rho=(r_1,\ldots,r_\ell)\vDash n$ run over all ordered compositions of
$n$, and let $s_j=r_1+\cdots+r_j$.  Then
\begin{equation}\label{eq:composition-formula}
 \boxed{
 F(2^{-n})
 =2^{-\binom n2}
 \sum_{\rho\vDash n}
 \prod_{j=1}^{\ell}
 \frac1{(2^{s_j}-1)(r_j+1)!}.}
\end{equation}
\end{theorem}
This derivation explains the factors $2^{s_j}-1$: they are the closed forms of
geometric sums over the infinitely many binary scales in
\eqref{eq:G-product}.

\subsection{A lower-Hessenberg determinant}

The recurrence \eqref{eq:moment-recurrence} may be written
\begin{equation}\label{eq:an-recurrence}
 a_n=\sum_{k=0}^{n-1}\beta_{n,k}a_k,
 \qquad
 \beta_{n,k}=\frac1{(2^n-1)(n-k+1)!}.
\end{equation}
Define the $n\times n$ lower-Hessenberg matrix
\begin{equation}\label{eq:Hn-matrix}
 H_n=
 \begin{pmatrix}
 \beta_{1,0}&-1&0&\cdots&0\\
 \beta_{2,0}&\beta_{2,1}&-1&\ddots&\vdots\\
 \beta_{3,0}&\beta_{3,1}&\beta_{3,2}&\ddots&0\\
 \vdots&\vdots&\vdots&\ddots&-1\\
 \beta_{n,0}&\beta_{n,1}&\beta_{n,2}&\cdots&\beta_{n,n-1}
 \end{pmatrix}.
\end{equation}

\begin{proposition}[Determinant representation]\label{prop:determinant}
For $n\ge1$,
\begin{equation}\label{eq:determinant-formula}
 \boxed{
  F(2^{-n})=2^{-\binom n2}\det H_n.}
\end{equation}
\end{proposition}

\begin{proof}
Let $D_0=1$ and $D_n=\det H_n$.  Expanding the determinant along its last row,
with the $-1$ superdiagonal supplying the compensating cofactor signs, gives
\[
 D_n=\sum_{k=0}^{n-1}\beta_{n,k}D_k.
\]
This is exactly \eqref{eq:an-recurrence}, with the same initial value; hence
$D_n=a_n$.
\end{proof}

Expanding the determinant by paths through the nonzero Hessenberg band
recovers the composition sum \eqref{eq:composition-formula}.  Thus the
composition and determinant formulae are two standard expansions of the same
triangular resolvent.

\subsection{Cumulants, Bernoulli numbers, and Bell polynomials}

Taking the logarithm of \eqref{eq:G-product} gives another exact description.
The classical expansion
\begin{equation}\label{eq:log-E}
 \log\frac{e^z-1}{z}
 =\frac z2+\sum_{m=1}^{\infty}
 \frac{B_{2m}}{2m(2m)!}z^{2m}
\end{equation}
uses the Bernoulli numbers $B_{2m}$.  Summing over the dyadic scales yields
\begin{equation}\label{eq:log-G}
 \log G(z)
 =\frac z2+
 \sum_{m=1}^{\infty}
 \frac{B_{2m}}{2m(2^{2m}-1)}
 \frac{z^{2m}}{(2m)!}.
\end{equation}
Therefore the cumulants of $X$ are
\begin{equation}\label{eq:cumulants}
 \kappa_1=\frac12,
 \qquad
 \kappa_{2m}=\frac{B_{2m}}{2m(2^{2m}-1)},
 \qquad
 \kappa_{2m+1}=0\quad(m\ge1).
\end{equation}
The vanishing odd cumulants beyond the first are another expression of the
symmetry $X\overset d=1-X$.

Let $\Bcal_n$ denote the complete exponential Bell polynomial, normalized by
\begin{equation}\label{eq:Bell-definition}
 \exp\left(\sum_{j=1}^{\infty}x_j\frac{z^j}{j!}\right)
 =\sum_{n=0}^{\infty}\Bcal_n(x_1,\ldots,x_n)\frac{z^n}{n!}.
\end{equation}
Then $d_n=\Bcal_n(\kappa_1,\ldots,\kappa_n)$, and hence
\begin{theorem}[Bernoulli--Bell representation]\label{thm:Bell}
\begin{equation}\label{eq:Bell-dyadic}
 \boxed{
 F(2^{-n})
 =\frac{2^{-\binom n2}}{n!}
 \Bcal_n(\kappa_1,\ldots,\kappa_n),}
\end{equation}
where the $\kappa_j$ are given by \eqref{eq:cumulants}.
\end{theorem}
Equivalently, expanding the complete Bell polynomial gives the finite partition
sum
\begin{equation}\label{eq:cumulant-partition}
 F(2^{-n})
 =2^{-\binom n2}
 \sum_{\substack{\nu_1,\ldots,\nu_n\ge0\\
                    \nu_1+2\nu_2+\cdots+n\nu_n=n}}
 \prod_{j=1}^{n}
 \frac1{\nu_j!}
 \left(\frac{\kappa_j}{j!}\right)^{\nu_j}.
\end{equation}
Only $\kappa_1$ and the even cumulants are nonzero, so the apparent
$n$-variable sum is substantially sparse.

For example,
\[
 d_2=\kappa_1^2+\kappa_2=\frac5{18},
 \qquad
 d_4=\kappa_1^4+6\kappa_1^2\kappa_2+3\kappa_2^2+\kappa_4
 =\frac{143}{1350}.
\]
Consequently,
\[
 F\!\left(\frac14\right)=\frac5{72},
 \qquad
 F\!\left(\frac1{16}\right)=\frac{143}{2073600}.
\]
The Bell-polynomial formula replaces Thue--Morse cancellation by cumulant
addition: the dyadic scales contribute independently to $\log G$.

\section{The discrete limit as a $q$-Toeplitz transform}

The discrete limit uses the same finite Gaussian coefficients as the exact
formula, but its inner sums are truncated at an $x$-dependent point rather
than taken over a complete Thue--Morse block.  Exact coefficient extraction is
therefore replaced by convergence of splines and regular summability.

\subsection{Shifted finite splines}

For $p\ge0$, $\eta\in\C$, and real $x$, define
\begin{equation}\label{eq:shifted-spline}
 S_{p,\eta}(x)=
 \frac{(-1)^p}{2^{\binom p2}p!}
 \sum_{0\le r<\lfloor2^px+1/2\rfloor}
 \eps_r\bigl(r-2^px+\eta\bigr)^p.
\end{equation}
An empty sum is zero.  The centered choice is $\eta=1/2$.
The formal development proves the uniform estimate
\begin{equation}\label{eq:centered-spline-error}
 |S_{p,1/2}(x)-\Fglobal(x)|\le2^{-p}
 \qquad(x\in\R),
\end{equation}
and, for fixed complex shift and $p\ge1$,
\begin{equation}\label{eq:shift-spline-error}
 |S_{p,\eta}(x)-S_{p,1/2}(x)|
 \le2^{1-p}e^{|\eta-1/2|}
 \qquad(x\ge0).
\end{equation}
On the negative half-line the relevant sums and the global extension vanish,
so the convergence statement extends without difficulty.

\subsection{The discrete $q$-sum}

Define
\begin{align}\label{eq:Dn-definition}
 D_n(\eta,x)
 &=\frac1{2^{n^2}\poch{1/2}{1/2}{n}}
 \sum_{k=0}^{n}
 \frac{\qbinom{n}{k}{1/2}}
 {4^{\binom k2}(n+k)!}
 \notag\\[-0.25em]
 &\quad\times
 \sum_{0\le r<\lfloor2^{n+k}x+1/2\rfloor}
 \eps_r\bigl(r-2^{n+k}x+\eta\bigr)^{n+k}.
\end{align}
This looks like a direct arbitrary-$x$ analogue of
\eqref{eq:centered-inverse-dyadic}.  The decisive simplification is an exact
reindexing.

\begin{theorem}[Exact Toeplitz decomposition]\label{thm:Toeplitz-decomp}
Put
\begin{equation}\label{eq:wnj-definition}
 w_{n,j}=
 \frac{(-1)^j\qbinom{n}{j}{1/2}
 2^{-\binom{j+1}{2}}}
 {\poch{1/2}{1/2}{n}}
 \qquad(0\le j\le n).
\end{equation}
Then
\begin{equation}\label{eq:Toeplitz-decomp}
 \boxed{
  D_n(\eta,x)=\sum_{j=0}^{n}w_{n,j}S_{2n-j,\eta}(x).}
\end{equation}
\end{theorem}

\begin{proof}
Set $j=n-k$ and $p=n+k=2n-j$ in \eqref{eq:Dn-definition}.  Replacing the
inner sum by the definition of $S_{p,\eta}$ leaves the power of $2$
\[
 \binom{2n-j}{2}-n^2-2\binom{n-j}{2}
 =-\binom{j+1}{2},
\]
and the sign $(-1)^{2n-j}=(-1)^j$.  Symmetry of the Gaussian binomial
coefficient gives
$\qbinom{n}{n-j}{1/2}=\qbinom{n}{j}{1/2}$.
\end{proof}

The indices $2n-j$ lie between $n$ and $2n$.  Thus every spline entering the
$n$th row is already at resolution at least $n$.

\subsection{One polynomial explains the entire row}

Define the row-generating polynomial
\begin{equation}\label{eq:Wn-definition}
  W_n(z)=\sum_{j=0}^{n}w_{n,j}z^j.
\end{equation}
The finite $q$-binomial theorem gives the compact product
\begin{theorem}[Toeplitz row polynomial]\label{thm:row-polynomial}
\begin{equation}\label{eq:row-polynomial}
 \boxed{
 W_n(z)=\frac{\poch{z/2}{1/2}{n}}
              {\poch{1/2}{1/2}{n}}.}
\end{equation}
\end{theorem}

Three important facts are immediate.

First, the row mass is one:
\begin{equation}\label{eq:row-mass}
 \sum_{j=0}^{n}w_{n,j}=W_n(1)=1.
\end{equation}
Second, the signs alternate, so the exact total variation is
\begin{equation}\label{eq:row-TV}
 V_n:=\sum_{j=0}^{n}|w_{n,j}|
 =W_n(-1)
 =\frac{\poch{-1/2}{1/2}{n}}
        {\poch{1/2}{1/2}{n}}.
\end{equation}
The sequence $V_n$ is bounded; the formal proof uses the convenient uniform
bound $V_n\le16$.

Third, differentiating at $z=1$ gives the first signed moment of the row:
\begin{equation}\label{eq:row-first-moment}
 \sum_{j=0}^{n}j\,w_{n,j}
 =W_n'(1)
 =-\sum_{\ell=1}^{n}\frac1{2^\ell-1}.
\end{equation}
Higher factorial moments follow by further logarithmic differentiation of
\eqref{eq:row-polynomial}.  Thus the entire Toeplitz stencil is encoded by one
finite $q$-Pochhammer ratio.

The weights also admit the symmetric form
\begin{equation}\label{eq:weights-symmetric}
 w_{n,j}=
 \frac{(-1)^j2^{-\binom{j+1}{2}}}
 {\poch{1/2}{1/2}{j}\poch{1/2}{1/2}{n-j}}.
\end{equation}
For fixed $j$ they converge to
\begin{equation}\label{eq:limit-weights}
 w_{\infty,j}=
 \frac{(-1)^j2^{-\binom{j+1}{2}}}
 {\poch{1/2}{1/2}{j}\poch{1/2}{1/2}{\infty}},
\end{equation}
whose generating function is
\begin{equation}\label{eq:limit-filter}
 W_\infty(z)=
 \frac{\poch{z/2}{1/2}{\infty}}
 {\poch{1/2}{1/2}{\infty}}.
\end{equation}
This limiting object is a stationary signed $q$-filter of unit mass and finite
variation.

\subsection{A quantitative discrete-limit estimate}

The exact Toeplitz decomposition makes an explicit error bound immediate.

\begin{theorem}[Uniform rate for fixed shift]\label{thm:Dn-rate}
For $n\ge1$, $x\ge0$, and $\eta\in\C$,
\begin{equation}\label{eq:Dn-rate}
 \boxed{
 |D_n(\eta,x)-\Fglobal(x)|
 \le V_n\left(2^{-n}+2^{1-n}e^{|\eta-1/2|}\right),}
\end{equation}
where $V_n$ is given by \eqref{eq:row-TV}.  In particular,
\begin{equation}\label{eq:Dn-rate-16}
 |D_n(\eta,x)-\Fglobal(x)|
 \le16\left(2^{-n}+2^{1-n}e^{|\eta-1/2|}\right).
\end{equation}
\end{theorem}

\begin{proof}
By \eqref{eq:row-mass} and \cref{thm:Toeplitz-decomp},
\[
 D_n(\eta,x)-\Fglobal(x)
 =\sum_{j=0}^{n}w_{n,j}
  \bigl(S_{2n-j,\eta}(x)-\Fglobal(x)\bigr).
\]
For $p=2n-j\ge n$, combine
\eqref{eq:centered-spline-error} and \eqref{eq:shift-spline-error}:
\[
 |S_{p,\eta}(x)-\Fglobal(x)|
 \le2^{-p}+2^{1-p}e^{|\eta-1/2|}
 \le2^{-n}+2^{1-n}e^{|\eta-1/2|}.
\]
Sum against $|w_{n,j}|$ and use \eqref{eq:row-TV}.
\end{proof}

For the centered shift one should not insert $\eta=1/2$ into the deliberately
coarse shift estimate.  The shift difference is then exactly zero, and the
sharper bound is
\begin{equation}\label{eq:Dn-centered-rate}
 |D_n(1/2,x)-\Fglobal(x)|\le V_n2^{-n}\le16\,2^{-n}.
\end{equation}

This argument shows precisely why the translation parameter disappears in the
limit.  It is not cancelled algebraically as in the finite exact formula;
rather, every fixed shift perturbs a degree-$p$ spline by $O(2^{-p})$, and the
Toeplitz rows have uniformly bounded total variation.

\subsection{Exact formula versus discrete limit}

The two constructions may now be compared without ambiguity.

\begin{center}
\renewcommand{\arraystretch}{1.25}
\begin{tabularx}{0.94\textwidth}{>{\raggedright\arraybackslash}p{0.23\textwidth}XX}
\toprule
 & Exact dyadic/Appell identity & Discrete limit \\
\midrule
Inner Thue--Morse sum
 & Complete block $0\le r<2^k$; exact polynomial annihilation
 & Prefix $0\le r<\lfloor2^{n+k}x+1/2\rfloor$; spline approximation \\
Outer Gaussian stencil
 & Extracts one coefficient exactly
 & Averages spline degrees $n,\ldots,2n$ with unit mass \\
Translation parameter
 & Cancels identically in the full numerator polynomial
 & Produces an $O(2^{-p})$ perturbation that vanishes in the limit \\
Passage to arbitrary $x$
 & Sum exact blocks over the binary $1$-digits of $x$
 & Let spline degree tend to infinity through a regular Toeplitz transform \\
\bottomrule
\end{tabularx}
\end{center}

The formulae share the same algebraic stencil, but the reason they are valid is
different.  Conflating coefficient extraction with Toeplitz convergence hides
the most interesting part of the story.

\section{Integral-transform representations}

The previous formulae read the entire function $G$ through its Taylor
coefficients.  Fourier and Laplace inversion read the same product globally.
They provide useful alternative representations at arbitrary, not necessarily
dyadic, arguments.

\subsection{Fourier--sinc inversion}

Let $Y=2X-1$.  Its density is Rvachev's up-function on $[-1,1]$.  With the
Fourier convention
\[
 \widehat f(\xi)=\int_{\R}f(t)e^{-2\pi i\xi t}\dd t,
\]
independence gives
\begin{equation}\label{eq:up-Fourier}
 \widehat{\Up}(\xi)
 =\prod_{j=0}^{\infty}
 \sinc\!\left(\frac{\pi\xi}{2^j}\right),
 \qquad
 \sinc u=\begin{cases}\sin u/u,&u\ne0,\\1,&u=0.\end{cases}
\end{equation}
Since $F(x)=\Prob(Y\le2x-1)$ and $\Up$ is even,
Fourier inversion followed by one integration gives, for $0<x<1$,
\begin{equation}\label{eq:F-Fourier}
 \boxed{
 F(x)=\frac12+\frac1\pi\int_0^{\infty}
 \frac{\sin\bigl(2\pi\xi(2x-1)\bigr)}{\xi}
 \prod_{j=0}^{\infty}
 \sinc\!\left(\frac{\pi\xi}{2^j}\right)\dd\xi.}
\end{equation}
The rapid decay of the infinite sinc product makes this a genuine convergent
integral, not merely a distributional formula.

The dyadic exact formulae are invisible term by term in
\eqref{eq:F-Fourier}; their rationality results from highly nonlocal
cancellation in the integral.  This contrast is useful: the Fourier form is
well adapted to smoothness and global shape, while the Appell and $q$-forms are
adapted to arithmetic at dyadic points.

\subsection{Bromwich inversion of the moment product}

The Laplace transform of $X$ is
\begin{equation}\label{eq:LX}
 L_X(s)=\E e^{-sX}=G(-s)
 =\prod_{j=1}^{\infty}
 \frac{1-e^{-s/2^j}}{s/2^j}.
\end{equation}
For $\sigma>0$ and $0<x<1$, inversion of the transform of the distribution
function gives
\begin{equation}\label{eq:Bromwich}
 \boxed{
 F(x)=\frac1{2\pi i}
 \lim_{T\to\infty}
 \int_{\sigma-iT}^{\sigma+iT}
 \frac{e^{sx}}{s}
 \prod_{j=1}^{\infty}
 \frac{1-e^{-s/2^j}}{s/2^j}\dd s.}
\end{equation}
The same product therefore supports three complementary readings:
\begin{itemize}
\item coefficient extraction gives moments and exact dyadic values;
\item logarithmic expansion gives cumulants and Bell-polynomial formulae;
\item contour inversion gives $F(x)$ directly at an arbitrary argument.
\end{itemize}

\section{Computational consequences}

\subsection{Exact dyadic evaluation}

For exact rational computation at $x=M/2^N$, the shortest route is usually:

\begin{enumerate}
\item compute $d_0,\ldots,d_N$ from the triangular recurrence
  \eqref{eq:moment-recurrence};
\item form $A_m(\theta)$ by Horner evaluation or by the binomial formula
  \eqref{eq:An-definition};
\item sum only over the $1$-bits in \eqref{eq:dyadic-digit-Appell}.
\end{enumerate}

This uses $O(N^2)$ rational arithmetic in a straightforward implementation and
avoids the exponentially long inner Thue--Morse sums that appear in the
closed $q$-formula.  The latter remain invaluable as structural identities and
formal certificates, but are not generally the numerically cheapest way to
obtain a rational value.

For inverse powers $x=2^{-N}$, only one bit is active, so
\eqref{eq:dyadic-digit-Appell} collapses at once to
\eqref{eq:inverse-dyadic-moment}.  The recurrence, composition sum,
determinant, Bell formula, and Gaussian-binomial formula then provide five
exactly equivalent ways to compute the same rational number.

\subsection{Arbitrary arguments}

For a non-dyadic $x$, the digit--Appell series
\eqref{eq:digit-Appell} is attractive because:
\begin{itemize}
\item zero bits cost nothing;
\item each active term needs one polynomial value $A_m(\theta_m)$;
\item the uniform remainder bound \eqref{eq:digit-tail-bound} is
  superexponentially small;
\item no cancellation inside a $2^k$-term Thue--Morse block is required.
\end{itemize}

The centered nested $q$-series \eqref{eq:centered-digit-q} is mathematically
equivalent term by term, but may be badly conditioned in ordinary floating
point because large powers cancel.  The raw translation freedom
\eqref{eq:raw-digit-q} can improve conditioning, although exact rational or
ball arithmetic is preferable when evaluating the nested sums directly.

The discrete limit \eqref{eq:Dn-definition} supplies an independent numerical
route and a global approximation to $\Fglobal$.  Its error is controlled by
\eqref{eq:Dn-rate}; the centered shift $\eta=1/2$ removes the shift-error term
entirely, giving the sharper estimate \eqref{eq:Dn-centered-rate}, and is
normally the best-conditioned canonical choice.

\section{Synthesis}

The apparent diversity of the formulae is reduced by separating three layers.

\paragraph{The moment layer.}
The random series \eqref{eq:X-definition} gives the entire moment product
\eqref{eq:G-product}.  Its coefficients are $d_n$, its logarithm gives the
cumulants \eqref{eq:cumulants}, and its dilation gives the recurrence
\eqref{eq:moment-recurrence}.

\paragraph{The finite polynomial layer.}
The reduction block is the rescaled Appell polynomial
\[
 R_n(2^{-n}\theta)=2^{-\binom n2}A_n(\theta)/n!.
\]
The finite Thue--Morse operator removes $k$ degrees, and the geometric
Gaussian-binomial stencil removes the remaining lower-degree terms.  This
produces the centered and translated finite $q$-formulae for the same Appell
polynomial.

\paragraph{The passage to arbitrary arguments.}
There are two distinct routes.
\begin{enumerate}
\item Follow the binary digits of $x$ and sum exact Appell blocks.  This yields
  \eqref{eq:digit-Appell}; substitution of the finite $q$-identity yields
  \eqref{eq:centered-digit-q} and \eqref{eq:raw-digit-q}.  Dyadic rationals are
  exactly the terminating cases.
\item Truncate Thue--Morse power sums at the spatial point $x$, interpret them
  as splines, and apply a bounded-variation Toeplitz row.  The row polynomial
  \eqref{eq:row-polynomial} makes the $q$-Pochhammer structure explicit, while
  \eqref{eq:Dn-rate} proves convergence quantitatively.
\end{enumerate}

The central identities can be compressed into the following chain:
\begin{equation}\label{eq:final-chain}
\boxed{
\begin{aligned}
 F(2^{-n})
 &=2^{-\binom n2}\frac{A_n(0)}{n!},\\[0.25em]
 R_n(2^{-n}\theta)
 &=2^{-\binom n2}\frac{A_n(\theta)}{n!},\\[0.25em]
 F(x)
 &=\sum_{m\ge1}b_m\eps_{a_{m-1}}
   2^{-\binom m2}\frac{A_m(\theta_m)}{m!},\\[0.25em]
 \frac{A_m(\theta)}{m!}
 &=\text{one finite Gaussian-binomial/Thue--Morse block},\\[0.25em]
 D_n(\eta,x)
 &=\sum_{j=0}^nw_{n,j}S_{2n-j,\eta}(x),
 \qquad
 \sum_{j=0}^nw_{n,j}z^j
 =\frac{\poch{z/2}{1/2}{n}}{\poch{1/2}{1/2}{n}}.
\end{aligned}}
\end{equation}
The first four lines are exact finite algebra assembled along the binary
expansion; the last line is regular summability of convergent spline
approximants.  Together they explain both the arithmetic exactness at dyadic
rationals and the analytic convergence at arbitrary arguments.

\appendix

\section{A compact proof of the geometric simplex sum}

For positive integers $r_1,\ldots,r_\ell$, put
$t_h=r_h+\cdots+r_\ell$.  Every increasing sequence
$1\le j_1<\cdots<j_\ell$ has a unique gap representation
\[
 j_i=i+g_1+\cdots+g_i,
 \qquad g_h\in\N\cup\{0\}.
\]
Therefore
\begin{align*}
 \sum_{j_1<\cdots<j_\ell}2^{-\sum_i j_ir_i}
 &=2^{-\sum_i ir_i}
   \prod_{h=1}^{\ell}
   \sum_{g_h=0}^{\infty}2^{-g_ht_h}\\
 &=2^{-\sum_i ir_i}
   \prod_{h=1}^{\ell}\frac1{1-2^{-t_h}}.
\end{align*}
But $\sum_ht_h=\sum_iir_i$, so
\[
 2^{-\sum_iir_i}\prod_h\frac{2^{t_h}}{2^{t_h}-1}
 =\prod_h\frac1{2^{t_h}-1},
\]
which is \eqref{eq:geometric-simplex}.

\section{First moments and inverse-dyadic values}

The recurrence \eqref{eq:moment-recurrence} gives:
\begin{center}
\renewcommand{\arraystretch}{1.18}
\begin{tabular}{ccl}
\toprule
$n$ & $d_n=\E X^n$ & $F(2^{-n})=2^{-\binom n2}d_n/n!$\\
\midrule
0 & $1$ & $1$\\
1 & $1/2$ & $1/2$\\
2 & $5/18$ & $5/72$\\
3 & $1/6$ & $1/288$\\
4 & $143/1350$ & $143/2073600$\\
\bottomrule
\end{tabular}
\end{center}
The $n=0$ row uses the standard convention $2^0=1$ and $0!=1$.

\section{Map to the formal development}

The principal formal ingredients used here are organized in the following
source files in the \texttt{Analysis/FabiusFunction} directory of the
\texttt{ProveIt} repository:

\begin{center}
\renewcommand{\arraystretch}{1.16}
\begin{tabularx}{0.96\textwidth}{>{\ttfamily\raggedright\arraybackslash}p{0.39\textwidth}X}
\toprule
\normalfont File & Mathematical role\\
\midrule
HalfQBinomial.lean
 & finite $q$-Pochhammer identities, Gaussian binomial coefficients at
   $q=1/2$, and the finite $q$-binomial annihilator\\
FabiusQBinomialFormula.lean
 & centered exact formulae for inverse and general dyadic arguments\\
FabiusRawQBinomialFormula.lean
 & raw-coordinate formula with arbitrary common translation\\
FabiusQBinomialTaylor.lean
 & polynomial identity equating the $q$-binomial block with the Taylor
   reduction polynomial\\
TaylorReduction.lean
 & exact reduction polynomial $R_n$ and dyadic-shell Taylor identities\\
FabiusBinaryReductionSeries.lean
 & finite telescoping identity, residual estimate, and global binary series\\
FabiusGlobalQBinomialSeries.lean
 & termwise substitution of the finite $q$-block into the binary series\\
FabiusUniformSpline.lean
 & centered finite splines and uniform convergence\\
FabiusComplexShiftSpline.lean
 & complex translation of the spline and quantitative shift stability\\
FabiusDiscreteLimitToeplitz.lean
 & Toeplitz weights, row mass, and total-variation bound\\
FabiusDiscreteLimitIntegration.lean
 & exact reindexing of the discrete formula and convergence to the global
   Fabius extension\\
\bottomrule
\end{tabularx}
\end{center}

\begin{thebibliography}{9}

\bibitem{Companion}
V.~Reshetnikov,
\emph{The Fabius Function and Rvachev's up-Function},
companion article in the \texttt{ProveIt} repository, 2026.

\bibitem{ProveIt}
V.~Reshetnikov,
\emph{ProveIt: formal developments for the Fabius function},
\url{https://github.com/VladimirReshetnikov/ProveIt/tree/main/Analysis/FabiusFunction}.

\bibitem{GasperRahman}
G.~Gasper and M.~Rahman,
\emph{Basic Hypergeometric Series}, 2nd ed.,
Cambridge University Press, 2004.

\bibitem{Comtet}
L.~Comtet,
\emph{Advanced Combinatorics},
D.~Reidel, 1974.

\bibitem{Hardy}
G.~H.~Hardy,
\emph{Divergent Series},
Oxford University Press, 1949.

\end{thebibliography}

\end{document}
